% !TEX program = xelatex
\documentclass[11pt,reqno]{amsart}

\usepackage{amsmath,amssymb,amsthm,mathtools}
\usepackage{booktabs}
\usepackage{enumitem}
\usepackage{float}
\usepackage{aliascnt}
\usepackage[no-math]{fontspec}
\usepackage{xeCJK}
\setCJKmainfont{Noto Serif CJK SC}
\setCJKmonofont{Noto Sans Mono CJK SC}
\usepackage[hidelinks]{hyperref}

\setlength{\textwidth}{6.15in}
\setlength{\textheight}{8.65in}
\setlength{\oddsidemargin}{0.18in}
\setlength{\evensidemargin}{0.18in}
\setlength{\topmargin}{-0.15in}

\newtheorem{theorem}{Theorem}[section]
\newaliascnt{lemma}{theorem}
\newtheorem{lemma}[lemma]{Lemma}
\aliascntresetthe{lemma}
\newaliascnt{proposition}{theorem}
\newtheorem{proposition}[proposition]{Proposition}
\aliascntresetthe{proposition}
\newaliascnt{corollary}{theorem}
\newtheorem{corollary}[corollary]{Corollary}
\aliascntresetthe{corollary}
\theoremstyle{definition}
\newaliascnt{definition}{theorem}
\newtheorem{definition}[definition]{Definition}
\aliascntresetthe{definition}
\theoremstyle{remark}
\newaliascnt{remark}{theorem}
\newtheorem{remark}[remark]{Remark}
\aliascntresetthe{remark}

\usepackage[nameinlink,capitalise,noabbrev]{cleveref}

\newcommand{\Sz}{\operatorname{Sz}}
\newcommand{\dist}{\operatorname{dist}}
\newcommand{\etaG}{\eta}
\newcommand{\N}{\mathbb{N}}
\newcommand{\cN}[1]{N[#1]}
\newcommand{\pairc}{\pi}
\newcommand{\Knt}[2]{K_{#1}^{#2}}

\title{The Strengthened Szeged--Wiener Conjecture}
\author{Zhijie He}
\date{}

\subjclass[2020]{05C12, 05C35, 05C85}
\keywords{Szeged index, Wiener index, distance in graphs, 2-connected graph}

\begin{document}
\begin{abstract}
Bonamy, Knor, Lu\v{z}ar, Pinlou and \v{S}krekovski conjectured in 2017 that every
$2$-connected graph $G$ of order $n\ge 10$, other than $K_n$, $K_n^2$ and
$K_n^{n-2}$, satisfies
\[
  \Sz(G)-W(G)\ge 2n.
\]
We prove the conjecture.  The argument completes the dominated-vertex induction
introduced in their work.  Complete deletions are handled by an exact formula for
$K_n^t$; the two exceptional $2$-connected deletions are treated by direct
true-twin counts; and non-$2$-connected deletions are reduced to a short list of
small block-pieces.  The finite part of the proof is confined to graphs of order at
most ten and is verified by an exact, reproducible certificate supplied with the
paper.  No floating-point computation or graph-isomorphism reduction is used in the
critical order-ten check.
\end{abstract}

\maketitle

\section{Introduction}

For a connected graph $G$, the Wiener index is
\[
  W(G)=\sum_{\{a,b\}\subseteq V(G)} d_G(a,b),
\]
where $d_G$ is graph distance.  If $e=uv$ is an edge, let $n_e(u)$ be the
number of vertices strictly closer to $u$ than to $v$, and define $n_e(v)$
analogously.  The Szeged index is
\[
  \Sz(G)=\sum_{uv\in E(G)} n_{uv}(u)n_{uv}(v).
\]
We write
\[
  \eta(G)=\Sz(G)-W(G).
\]

Bonamy, Knor, Lu\v{z}ar, Pinlou and \v{S}krekovski proved that every
$2$-connected non-complete graph of order $n$ satisfies
\(\eta(G)\ge 2n-6\) \cite[Theorem~2]{BKLPS}.  Their stronger
Theorem~4, together with Lemma~11, identifies the equality graphs as
$K_n^2$ and $K_n^{n-2}$ and gives the bound $2n-5$ after excluding them
\cite[Theorem~4 and Lemma~11]{BKLPS}.  They then proposed the following
strengthening \cite[Conjecture~5]{BKLPS}.

For $1\le t\le n-1$, let $K_n^t$ denote the graph obtained from $K_{n-1}$ by
adding one vertex adjacent to exactly $t$ vertices of the old clique.  Thus
$K_n^{n-1}=K_n$.

\begin{theorem}\label{thm:main}
Let $G$ be a $2$-connected graph of order $n\ge 10$.  If
\[
  G\not\cong K_n,\qquad G\not\cong K_n^2,\qquad
  G\not\cong K_n^{n-2},
\]
then
\[
  \eta(G)\ge 2n.
\]
\end{theorem}

The proof stays close to the mechanism of \cite{BKLPS}.  The useful point is
that a putative counterexample must contain a vertex whose closed neighbourhood
is contained in another closed neighbourhood.  Deleting such a vertex either
preserves $2$-connectivity, or forces a rigid block decomposition.  The former
case has two exceptional subcases not covered by a direct induction; the latter
case requires stronger low-order estimates than those used in
\cite[Theorem~4]{BKLPS}.  These are the two places where new input is needed.

The finite verification appears only after the structural reductions have been
proved.  We state the exact finite universes in \cref{sec:finite}; the
supplementary certificate merely evaluates the definitions on those universes.

\section{Good edges and dominated vertices}\label{sec:prelim}

We use the good-edge formulation of \cite{BKLPS}.  For an unordered pair
$\{a,b\}$, an edge $xy$ is \emph{good} for $\{a,b\}$ if, after possibly
interchanging $x$ and $y$,
\[
 d_G(a,x)<d_G(a,y),\qquad d_G(b,y)<d_G(b,x).
\]
Let $g_G(a,b)$ be the number of good edges for $\{a,b\}$ and set
\[
  \pairc_G(a,b)=g_G(a,b)-d_G(a,b).
\]

\begin{proposition}[Bonamy--Knor--Lu\v{z}ar--Pinlou--\v{S}krekovski]\label{prop:good}
For every connected graph $G$,
\[
  \eta(G)=\sum_{\{a,b\}\subseteq V(G)}\pairc_G(a,b),
\]
and every summand is non-negative.
\end{proposition}

\begin{proof}
This is \cite[Proposition~6]{BKLPS}.
\end{proof}

We shall also use the following three published lemmas.  They are quoted with
all hypotheses because the induction below depends on them exactly as stated.

\begin{lemma}[Dominated deletion]\label{lem:published7}
Let $G$ be $2$-connected.  Suppose that $u,v\in V(G)$ satisfy
$N[u]\subseteq N[v]$, and that $G-u$ is $2$-connected and non-complete.  Then
\[
  \eta(G)-\eta(G-u)\ge 2.
\]
\end{lemma}

\begin{proof}
This is \cite[Lemma~7]{BKLPS}.
\end{proof}

\begin{lemma}[Structure after a non-$2$-connected deletion]\label{lem:published9}
Let $G$ be a $2$-connected non-complete graph distinct from $K_4^2$.
Assume that at least one pair $u,v$ satisfies $N[u]\subseteq N[v]$, and that
for every such pair the graph $G-u$ is complete or is not $2$-connected.
Then, for every such pair, $G-u$ is not $2$-connected.  Moreover, the
dominator $v$ is the unique cut vertex of $G-u$, and no block of $G-u$ is
isomorphic to $K_t$, $K_t^2$, or $K_t^{t-2}$.
\end{lemma}

\begin{proof}
The first and last assertions are \cite[Lemma~9]{BKLPS}; the proof there also
shows explicitly that $v$ is the unique cut vertex of $G-u$.
\end{proof}

\begin{lemma}[No dominated pair]\label{lem:published10}
Let $G$ be a $2$-connected non-complete graph distinct from $C_5$.  If no two
vertices $u,v$ satisfy $N[u]\subseteq N[v]$, then
\[
  \eta(G)\ge 2|V(G)|.
\]
\end{lemma}

\begin{proof}
This is the final conclusion of \cite[Lemma~10]{BKLPS}.
\end{proof}

Two elementary observations will make the block argument transparent.

\begin{lemma}\label{lem:pairzero}
If $N[u]\subseteq N[v]$, then $uv\in E(G)$ and
\[
  \pairc_G(u,v)=0.
\]
\end{lemma}

\begin{proof}
The inclusion contains $u\in N[v]$, so $uv$ is an edge.  For $z=v$ the inequality $d(v,z)\le d(u,z)$ is immediate.  If
$z\notin\{u,v\}$, a shortest $u$--$z$ path begins with a neighbour of $u$;
that neighbour also lies in $N(v)$, and hence again $d(v,z)\le d(u,z)$.
Since $uv$ is an edge,
$d(u,z)\le d(v,z)+1$.

The edge $uv$ is good for $\{u,v\}$.  Suppose another edge $xy$ were good,
oriented so that $d(u,x)<d(u,y)$.  If $x=u$, then the strict inequality at
$v$ forces $y=v$, so the edge is $uv$.  Hence for another good edge we have
$x\ne u$, and the inequality above gives $d(u,x)\ge d(v,x)$.  Adjacency
gives $d(u,y)=d(u,x)+1$, while the edge $uv$ gives
\[
 d(v,y)\ge d(u,y)-1=d(u,x)\ge d(v,x),
\]
contradicting the reverse strict inequality required at $v$.  Thus
$g(u,v)=1=d(u,v)$.
\end{proof}

\begin{lemma}[Encoding a dominated extension]\label{lem:encoding}
Let $H=G-u$ and suppose $N_G[u]\subseteq N_G[v]$.  Then
\[
  N_G(u)=\{v\}\cup T
\]
for some $T\subseteq N_H(v)$.  If $G$ is $2$-connected, then $T\ne\varnothing$.
Conversely, adjoining to any graph $H$ a new vertex $u$ with neighbourhood
$\{v\}\cup T$, where $T\subseteq N_H(v)$, produces $N[u]\subseteq N[v]$.
If, in addition, $H$ is $2$-connected and $T\ne\varnothing$, then the resulting
graph is $2$-connected.
\end{lemma}

\begin{proof}
The closed-neighbourhood inclusion gives $uv\in E(G)$, and every other
neighbour of $u$ must be adjacent to $v$.  In a $2$-connected graph
$\deg(u)\ge2$, so $T$ is non-empty.  The converse inclusion is immediate.
For the final assertion, deleting $u$ leaves $H$; deleting a vertex of $H$
leaves the connected graph $H-x$, and at least one of the two distinct
neighbours $v$ and a vertex of $T$ remains adjacent to $u$.  Thus deletion of
any single vertex leaves a connected graph.
\end{proof}

\section{The exceptional \texorpdfstring{$2$}{2}-connected deletions}\label{sec:exceptional}

We first remove the case in which deleting the dominated vertex leaves a
complete graph.

\begin{proposition}\label{prop:knt}
For $2\le t\le n-2$,
\[
  \eta(K_n^t)=2(t-1)(n-t-1).
\]
\end{proposition}

\begin{proof}
Let $x$ be the added vertex, let $A=N(x)$ with $|A|=t$, and let
$B$ be the remaining $s=n-1-t$ vertices of the old clique.  Then
\[
 W(K_n^t)=\binom{n-1}{2}+t+2s.
\]
Every edge inside the old clique contributes $1$ to the Szeged index, except
that each of the $ts$ edges between $A$ and $B$ receives one additional unit
from $x$.  Each of the $t$ edges from $x$ to $A$ contributes $s+1$.  Hence
\[
 \Sz(K_n^t)=\binom{n-1}{2}+ts+t(s+1),
\]
and subtraction gives the formula.
\end{proof}

\begin{corollary}\label{cor:complete-deletion}
Let $n\ge10$.  If $G$ is $2$-connected, $N[u]\subseteq N[v]$, and $G-u$ is
complete, then either $G\cong K_n,K_n^2,K_n^{n-2}$ or $\eta(G)\ge2n$.
\end{corollary}

\begin{proof}
We have $G\cong K_n^t$.  The cases $t=1$ and $t=n-1$ are, respectively,
not $2$-connected and complete.  After excluding $t=2,n-2$, it remains to
consider $3\le t\le n-3$.  The expression in \cref{prop:knt} is concave in
$t$, so its minimum on this interval occurs at an endpoint and equals
$4n-16\ge2n$.
\end{proof}

The other obstruction to the direct induction is that $G-u$ may itself be one
of the two equality graphs identified by \cite[Theorem~4 and Lemma~11]{BKLPS}.
Their symmetry makes an exact calculation short.  We record the counting device once.

\begin{lemma}[True-twin quotient count]\label{lem:twin}
Suppose $V(G)$ is partitioned into non-empty classes
$V_1,\dots,V_r$ such that each $V_i$ is a clique of true twins (singletons are
allowed), and between two distinct classes all edges or no edges are present.
Let $m_i=|V_i|$, let $Q$ be the quotient graph, and write $D_{ij}$ for distance
in $Q$.  Then
\[
 W(G)=\sum_i\binom{m_i}{2}+\sum_{i<j}m_im_jD_{ij}.
\]
For an edge $ij$ of $Q$, put
\[
 A_{ij}=1+\sum_{h\notin\{i,j\}\,:\,D_{hi}<D_{hj}}m_h,
 \qquad
 B_{ij}=1+\sum_{h\notin\{i,j\}\,:\,D_{hj}<D_{hi}}m_h.
\]
Then
\[
 \Sz(G)=\sum_i\binom{m_i}{2}
       +\sum_{ij\in E(Q)}m_im_jA_{ij}B_{ij}.
\]
\end{lemma}

\begin{proof}
Distances are constant between two distinct classes and equal to the quotient
distance.  An internal edge of a true-twin clique has one vertex on each
strict side and all other vertices are equidistant, so its Szeged contribution
is $1$.  For a cross edge between $V_i$ and $V_j$, the other vertices in
$V_i\cup V_j$ are equidistant from its endpoints, while every third class
lies entirely on one strict side or is equidistant according to the quotient
distances.  This gives the displayed formulas.
\end{proof}

\subsection{Deletion to \texorpdfstring{$K_{n-1}^2$}{K(n-1,2)}}

Assume $H=G-u\cong K_{n-1}^2$.  Write the old clique as
\[
 Q=\{a,b\}\cup R,\qquad |R|=n-4,
\]
and let $x$ be the exceptional vertex, adjacent precisely to $a,b$ in $H$.
Set
\[
 \varepsilon=\mathbf 1_{ux\in E(G)},\qquad
 k=|N(u)\cap\{a,b\}|,\qquad
 \ell=|N(u)\cap R|.
\]

\begin{lemma}\label{lem:k2patterns}
If $G$ is $2$-connected and $N[u]\subseteq N[v]$ for some $v\in V(H)$,
then $(\varepsilon,k)$ is one of
\[
 (0,0),(0,1),(0,2),(1,1),(1,2),
\]
with the ranges for $\ell$ shown in \cref{tab:k2}.
\end{lemma}

\begin{proof}
If $v=x$, then $u$ is adjacent to $x$, to no vertex of $R$, and to at least
one of $a,b$.  If $v\in\{a,b\}$, that vertex is dominating in $H$, so the
only condition is that it be adjacent to $u$, hence $k\ge1$.  If $v\in R$,
then $x\notin N[v]$, so $\varepsilon=0$.  These three possibilities give
exactly the five rows.  The lower bounds on $\ell$ are precisely
$\deg(u)\ge2$; the upper bound is $|R|=n-4$.
\end{proof}

For an interior value $0<\ell<n-4$, split $\{a,b\}$ into the selected and
unselected classes $S,T$ of sizes $k,2-k$, and split $R$ into the selected
and unselected classes $L,M$ of sizes $\ell,n-4-\ell$.  The classes
$S,T,L,M$ form a clique; $x$ is adjacent exactly to $S\cup T$; and $u$ is
adjacent to $S\cup L$ and, when $\varepsilon=1$, to $x$.  Applying
\cref{lem:twin} gives the following exact expressions.  When $L$ or $M$ is
empty, the same endpoint values are obtained by omitting that class and
repeating the quotient count; thus no polynomial extrapolation is being used.

\begin{table}[H]
\centering
\caption{The case $G-u\cong K_{n-1}^2$.  The last column excludes the endpoint
that is isomorphic to $K_n^2$.}\label{tab:k2}
\small
\begin{tabular}{@{}c c l c@{}}
\toprule
$(\varepsilon,k)$ & admissible $\ell$ & $\eta(G)$ & non-exceptional minimum \\
\midrule
$(0,0)$ & $2\le\ell\le n-4$
& $-2\ell^2+(2n+1)\ell-5$ & $4n-11$ \\
$(0,1)$ & $1\le\ell\le n-4$
& $-2\ell^2+(2n-5)\ell+2n-8$ & $4n-15$ \\
$(0,2)$ & $0\le\ell\le n-4$
& $-2\ell^2+(2n-10)\ell+4n-14$ & $4n-14$ \\
$(1,1)$ & $0\le\ell\le n-4$
& $-3\ell^2+(3n-5)\ell+3n-10$ & $3n-10$ \\
$(1,2)$ & $0\le\ell\le n-4$
& $-3\ell^2+(3n-12)\ell+4n-16$ & $4n-16$ \\
\bottomrule
\end{tabular}
\end{table}

Each row is a concave quadratic in $\ell$.  The only exceptional endpoint is
$(\varepsilon,k,\ell)=(0,2,n-4)$, which is $K_n^2$ and has value $2n-6$.
Direct endpoint counts give, when $\ell=n-4$, respectively
\[
 9n-41,\quad 5n-20,\quad 2n-6,\quad 10n-38,\quad 4n-16,
\]
and when $\ell=0$ is admissible they give
\[
 4n-14,\quad 3n-10,\quad 4n-16
\]
for the last three relevant rows.  Hence the minima in \cref{tab:k2} follow.

For illustration, in the extremal row $(\varepsilon,k)=(1,1)$ and
$0<\ell<n-4$, the quotient classes have sizes
\[
  1,1,1,1,\ell,n-4-\ell
\]
for $\{u\},\{x\},S,T,L,M$.  Substitution in \cref{lem:twin} gives
\begin{align*}
 W(G)&=-\ell^2+n\ell-5\ell+6n-17
       +\binom{\ell}{2}+\binom{n-4-\ell}{2},\\
 \Sz(G)&=-4\ell^2+4n\ell-10\ell+9n-27
       +\binom{\ell}{2}+\binom{n-4-\ell}{2},
\end{align*}
so their difference is the fourth row of \cref{tab:k2}.  The other rows are
obtained by the same quotient calculation; the tables record the full
arithmetic needed later.

\begin{proposition}\label{prop:k2extension}
Let $n\ge10$.  Suppose $G$ is $2$-connected, $N[u]\subseteq N[v]$, and
$G-u\cong K_{n-1}^2$.  If $G\not\cong K_n^2$, then
\[
  \eta(G)\ge 3n-10\ge2n.
\]
\end{proposition}

\subsection{Deletion to \texorpdfstring{$K_{n-1}^{n-3}$}{K(n-1,n-3)}}

Now let $H=G-u\cong K_{n-1}^{n-3}$.  Write the old clique as
$Q=A\cup\{z\}$ with $|A|=n-3$, and let $x$ be adjacent to every vertex of
$A$ and not to $z$.  Set
\[
 \varepsilon=\mathbf 1_{ux\in E(G)},\qquad
 \delta=\mathbf 1_{uz\in E(G)},\qquad
 \ell=|N(u)\cap A|.
\]
If the dominator lies in $A$, it must be one of the $\ell$ selected vertices
and imposes no further restriction.  If it is $x$, then
$(\varepsilon,\delta)=(1,0)$; if it is $z$, then
$(\varepsilon,\delta)=(0,1)$.  Consequently all four pairs
$(0,0),(0,1),(1,0),(1,1)$ occur, with $\ell\ge2$ in the first row and
$\ell\ge1$ in the others.

Splitting $A$ into selected and unselected true-twin classes and applying
\cref{lem:twin} gives \cref{tab:kn3}.  The endpoint $\ell=n-3$ is checked
directly after deleting the empty unselected class.

\begin{table}[H]
\centering
\caption{The case $G-u\cong K_{n-1}^{n-3}$.  The exceptional endpoint in the
last row is $K_n^{n-2}$.}\label{tab:kn3}
\small
\begin{tabular}{@{}c c l c@{}}
\toprule
$(\varepsilon,\delta)$ & admissible $\ell$ & $\eta(G)$ & non-exceptional minimum \\
\midrule
$(0,0)$ & $2\le\ell\le n-3$
& $-2\ell^2+2n\ell-6$ & $4n-14$ \\
$(0,1)$ & $1\le\ell\le n-3$
& $-2\ell^2+(2n-5)\ell+3n-13$ & $4n-16$ \\
$(1,0)$ & $1\le\ell\le n-3$
& $-2\ell^2+(2n-5)\ell+3n-13$ & $4n-16$ \\
$(1,1)$ & $1\le\ell\le n-3$
& $-2\ell^2+(2n-12)\ell+8n-24$ & $4n-8$ \\
\bottomrule
\end{tabular}
\end{table}

At $\ell=n-3$ the four rows have values
\[
 6n-24,\quad 4n-16,\quad 4n-16,\quad 2n-6,
\]
respectively; the last graph is exactly $K_n^{n-2}$.  Concavity again gives
the final column of \cref{tab:kn3}.

\begin{proposition}\label{prop:kn3extension}
Let $n\ge10$.  Suppose $G$ is $2$-connected, $N[u]\subseteq N[v]$, and
$G-u\cong K_{n-1}^{n-3}$.  If $G\not\cong K_n^{n-2}$, then
\[
  \eta(G)\ge 4n-16\ge2n.
\]
\end{proposition}

\section{Block pieces}\label{sec:blocks}

Suppose now that $G$ is $2$-connected, $N[u]\subseteq N[v]$, and $G-u$ is
not $2$-connected.  Under the hypotheses of \cref{lem:published9}, let
$C_1,\dots,C_k$ be the blocks of $G-u$.  The vertex $v$ is their unique
common cut vertex.  Put
\[
  G_i=G[V(C_i)\cup\{u\}],\qquad s_i=|V(G_i)|.
\]
Since $G$ is $2$-connected, $u$ has in every $C_i$ a neighbour
$w_i\ne v$; otherwise $v$ would be a cut vertex of $G$.  Hence each $G_i$ is
$2$-connected by \cref{lem:encoding}, applied to the block $C_i$.  By
\cref{lem:published9}, each $C_i$ is non-complete and is not $K_t^2$ or
$K_t^{t-2}$; in particular $|C_i|\ge4$ and $s_i\ge5$.

The proof of \cite[Theorem~4]{BKLPS} observes in precisely this
situation that every $G_i$ is a proper $2$-connected subgraph and is neither
complete nor isomorphic to $K_{s_i}^2$ or $K_{s_i}^{s_i-2}$.  For completeness,
we isolate the short reason.

\begin{lemma}\label{lem:pieces-nonspecial}
Each $G_i$ is $2$-connected and
\[
 G_i\not\cong K_{s_i},\qquad
 G_i\not\cong K_{s_i}^2,\qquad
 G_i\not\cong K_{s_i}^{s_i-2}.
\]
\end{lemma}

\begin{proof}
The $2$-connectivity was noted above.  If $G_i$ were complete, then $C_i$
would be complete.  In $K_{s_i}^2$, deleting the degree-$2$ vertex leaves
$K_{s_i-1}$; deleting one of the $s_i-3$ old non-universal clique vertices
leaves $K_{s_i-1}^2$; and deleting either universal vertex makes the degree-$2$
vertex pendant, so the deletion is not $2$-connected.  Hence a deletion that
remains $2$-connected is complete or $K_{s_i-1}^2$.

In $K_{s_i}^{s_i-2}$, deleting either of the two non-universal vertices leaves
$K_{s_i-1}$, whereas deleting a universal vertex leaves
$K_{s_i-1}^{s_i-3}$.  Thus, again, every $2$-connected deletion is complete or
one of the two families forbidden for $C_i$ by \cref{lem:published9}.  This
is the same case analysis used in \cite[proof of Theorem~4]{BKLPS}.
\end{proof}

We shall use the following small-piece lower bounds, proved in
\cref{sec:finite} by an exact finite certificate:
\begin{equation}\label{eq:Ltable}
\begin{array}{c|ccccc}
 s&5&6&7&8&9\\ \hline
 L(s)&10&8&11&14&17.
\end{array}
\end{equation}
Here $L(s)$ is a lower bound for every $s$-vertex $2$-connected graph of the
form in \cref{lem:encoding} whose deletion graph is a non-special block.
For $s\ge10$ occurring inside a minimal counterexample, the inductive bound
will be $L(s)=2s$.

\begin{lemma}[Block assembly]\label{lem:blockassembly}
With the notation above,
\[
  \eta(G)\ge \sum_{i=1}^k\eta(G_i)+k(k-1).
\]
\end{lemma}

\begin{proof}
The pieces meet only in $u$ and $v$.  A path between two vertices of a fixed
$G_i$ that leaves $G_i$ can enter another block only through $u$ or $v$;
since $uv$ is already an edge, replacing the outside detour by $uv$ never
increases the length.  Hence distances inside $G_i$ are the same in $G$.
Every edge of $G_i$ that is good for a pair inside $G_i$ therefore remains
good in $G$.  The only pair counted in more than one $G_i$ is $\{u,v\}$,
and its contribution is zero by \cref{lem:pairzero}.  Thus the contributions
from within the pieces sum to at least $\sum_i\eta(G_i)$.

For $i\ne j$, the vertices $w_i,w_j$ lie in distinct blocks of $G-u$ and so
are non-adjacent.  Both are adjacent to $u$ and, because
$N[u]\subseteq N[v]$, also to $v$.  The four edges
$w_i u,uw_j,w_i v,vw_j$ are good for $\{w_i,w_j\}$, while
$d(w_i,w_j)=2$.  Hence
\[
  \pairc_G(w_i,w_j)\ge2.
\]
Summing over the $\binom{k}{2}$ unordered pairs of blocks gives the extra
term $2\binom{k}{2}=k(k-1)$.
\end{proof}

Because $u$ and $v$ occur in every piece,
\[
  \sum_{i=1}^k s_i=n+2(k-1).
\]
Define $d(s)=L(s)-2s$.  From \eqref{eq:Ltable},
\[
 d(5)=0,\quad d(6)=-4,\quad d(7)=-3,\quad d(8)=-2,\quad d(9)=-1.
\]
For a piece of order at least ten inside a minimal counterexample we shall
have $d(s)\ge0$.  Therefore \cref{lem:blockassembly} yields
\begin{equation}\label{eq:deficit}
 \eta(G)\ge 2n+(k-1)(k+4)+\sum_{i=1}^k d(s_i).
\end{equation}

\begin{lemma}\label{lem:deficits}
If $k\ge3$, then the right-hand side of \eqref{eq:deficit} is at least
$2n+2$.  If $k=2$ and $n\ge12$, it is at least $2n$.  For $n=10,11$ the only
combinations not settled by \eqref{eq:deficit} are
\[
 n=10:\ (s_1,s_2)=(6,6),\qquad
 n=11:\ (s_1,s_2)=(6,7).
\]
\end{lemma}

\begin{proof}
Every deficit is at least $-4$.  For $k\ge3$, \eqref{eq:deficit} is at least
\[
 2n+(k-1)(k+4)-4k=2n+k^2-k-4\ge2n+2.
\]
If $k=2$ and $n\ge12$, then $s_1+s_2=n+2\ge14$ with each $s_i\ge5$.
Inspection of the five displayed deficits, together with $d(s)\ge0$ for
$s\ge10$, gives $d(s_1)+d(s_2)\ge-6$, which cancels the term
$(k-1)(k+4)=6$.

For $n=10$, $s_1+s_2=12$; the pair $(5,7)$ is already safe and only $(6,6)$
remains.  For $n=11$, $s_1+s_2=13$; $(5,8)$ is safe and only $(6,7)$ remains.
\end{proof}

\section{The finite certificate}\label{sec:finite}

This section states exactly what is verified computationally and why those
finite universes cover every case left by the preceding lemmas.  All values of
$W$, $\Sz$ and $\eta$ are computed as integers from their definitions.

The input small graphs are taken from the Atlas of Graphs of Read and Wilson
\cite{ReadWilson}.  The supplied snapshot was extracted with
\texttt{graph\_atlas\_g} from NetworkX~3.6.1 \cite{NetworkXAtlas}; the
NetworkX documentation identifies this routine as the list of all graphs on at
most seven vertices named in the Atlas.  The certificate contains the fixed
snapshot and its SHA256 digest.  Regeneration is optional: the C++ verifier
runs directly from the snapshot.

For a graph $C$, a vertex $v\in V(C)$ and a non-empty set
$T\subseteq N_C(v)$, write
\[
  P(C,v,T)
\]
for the graph obtained by adjoining a new vertex $u$ with
$N(u)=\{v\}\cup T$.  By \cref{lem:encoding}, every block-piece in the proof
has this form.

\begin{proposition}[Small-piece certificate]\label{prop:pieces}
Let $5\le s\le9$.  Let $C$ be a $2$-connected graph of order $s-1$ that
is not isomorphic to $K_{s-1}$, $K_{s-1}^2$, or
$K_{s-1}^{s-3}$.  For every $v\in V(C)$ and every non-empty
$T\subseteq N_C(v)$,
\[
  \eta(P(C,v,T))\ge L(s),
\]
where the exact minima are
\[
 L(5)=10,\quad L(6)=8,\quad L(7)=11,\quad L(8)=14,\quad L(9)=17.
\]
\end{proposition}

\begin{proof}
For $5\le s\le8$, the deletion block $C$ has $4\le|C|\le7$.  The certificate
runs over every $2$-connected Atlas representative $C$ of that order other
than $K_t,K_t^2,K_t^{t-2}$, every choice of $v$, and every non-empty
$T\subseteq N_C(v)$.  The numbers of generated parameter choices and the
resulting minima are
\[
\begin{array}{c|rrrr}
 s&5&6&7&8\\ \hline
 \#P(C,v,T)&12&209&3258&48865\\
 \min\eta&10&8&11&14.
\end{array}
\]

For $s=9$, let $C$ have order eight.  If $\eta(C)\ge15$, then
\cref{lem:published7} gives $\eta(P(C,v,T))\ge17$.  It remains only to
consider $2$-connected, non-special $C$ with $\eta(C)\le14$.  Every
$2$-connected eight-vertex graph becomes connected after deletion of any
vertex, so it is obtained by adding one non-isolated vertex to a connected
seven-vertex Atlas representative.  The certificate generates all such
extensions; 65 generated occurrences survive the stated filters.  Every
dominated extension of these occurrences is then checked, and the minimum is
$17$.  Repeated isomorphic occurrences are deliberately retained.
\end{proof}

We next isolate the two small gluings left by \cref{lem:deficits}.

\begin{proposition}[Glue certificate]\label{prop:glue}
In the block-decomposition setting of \cref{sec:blocks}, suppose $k=2$.
If $n=10$ and $(s_1,s_2)=(6,6)$, then $\eta(G)\ge20$; if $n=11$ and
$(s_1,s_2)=(6,7)$, then $\eta(G)\ge23$.
\end{proposition}

\begin{proof}
In the first case \cref{lem:blockassembly} already gives $\eta(G)\ge20$
whenever $\eta(G_1)+\eta(G_2)\ge18$.  Since $L(6)=8$, only $6$-vertex pieces
of values $8$ or $9$ need be retained, and only pairs with sum below $18$
need be glued.  The piece generation in \cref{prop:pieces} yields 15 such
piece occurrences; there are 105 relevant unordered gluing configurations.
Their exact minimum is $20$.

For $(6,7)$, the lower bounds $L(6)=8$ and $L(7)=11$ show that only pieces
attaining these two minima can possibly give a value below $22$.  There are
10 and 16 such occurrences, respectively, producing 160 gluings.  Their
minimum is $23$.
\end{proof}

The last finite statement supplies the induction base at order ten when a
dominated deletion remains $2$-connected.

\begin{proposition}[Order-ten base certificate]\label{prop:base10}
Let $G$ be a $2$-connected graph of order ten, distinct from
$K_{10},K_{10}^2,K_{10}^8$.  Then $\eta(G)\ge20$.
\end{proposition}

\begin{proof}
Suppose instead that $\eta(G)\le19$.  By \cref{lem:published10}, $G$ has a
dominated pair $u,v$.  If some dominated deletion is complete, or is one of the two exceptional
$2$-connected graphs, then
\cref{cor:complete-deletion,prop:k2extension,prop:kn3extension} already give
a contradiction.  If no dominated deletion is $2$-connected, then every one
is non-$2$-connected and \cref{lem:published9,lem:deficits,prop:glue} settle
the graph.  Hence we may choose a dominated pair such that $H=G-u$ is
$2$-connected, non-complete and non-special.  By \cref{lem:published7},
\[
  \eta(H)\le17.
\]

Since $\eta(H)\le17<18$, the contrapositive of
\cref{lem:published10} gives a dominated pair in the nine-vertex graph $H$.
If all dominated deletions of $H$ were complete or non-$2$-connected, then
\cref{lem:published9} and the block estimate would force the only possible
piece sizes to be $(5,6)$ and hence
\[
 \eta(H)\ge L(5)+L(6)+2=20,
\]
a contradiction.  Thus $H$ has a dominated deletion $F=H-a$ that is
$2$-connected or complete.  If $F$ is complete, then $H\cong K_9^t$; by
\cref{prop:knt}, every non-special such graph has $\eta(H)\ge20$, again a
contradiction.  Hence $F$ is $2$-connected and non-complete, and
\cref{lem:published7} gives
\[
  \eta(F)\le15.
\]

The finite universe can now be generated without an isomorphism search.
Start with all 853 connected seven-vertex Atlas representatives and add one
non-isolated vertex in every possible way; retain the $2$-connected graphs
$F$ with $\eta(F)\le15$.  This produces 97 labelled occurrences.  From each,
generate every dominated extension from \cref{lem:encoding}; retain the
non-special graphs $H$ with $\eta(H)\le17$.  This produces 681 occurrences.
Finally generate every dominated extension of each retained $H$, excluding
only the three special order-ten graphs.  At this last stage no
$2$-connectivity filter and no isomorphism deduplication are used, so the
300311 generated graphs form a superset of all remaining candidates.  Their
minimum value is exactly $20$, a contradiction.
\end{proof}

\begin{remark}[What the certificate checks]\label{rem:certificate}
The verifier performs all-pairs breadth-first search, sums the resulting
integer distances for $W$, and for every edge $uv$ counts directly the
vertices on the two strict distance sides to obtain the factor
$n_{uv}(u)n_{uv}(v)$.  A graph is tested for $2$-connectivity by deleting each
vertex in turn and checking connectivity.  The three exceptional families are
recognized by their degree sequences, which uniquely determine them:
\[
\begin{aligned}
 K_n &: (n-1)^n,\\
 K_n^2 &: 2,(n-2)^{n-3},(n-1)^2,\\
 K_n^{n-2} &: (n-2)^2,(n-1)^{n-2}.
\end{aligned}
\]
For $K_n^2$, the two vertices of degree $n-1$ are universal, so the
degree-$2$ vertex is forced to be adjacent exactly to them.  For
$K_n^{n-2}$, the $n-2$ vertices of degree $n-1$ are universal, forcing the
remaining two vertices to be non-adjacent.  Thus these degree tests introduce
no isomorphism ambiguity.  All arithmetic is integral.  The critical order-ten
run intentionally keeps
repeated labelled copies, so correctness does not depend on a graph-isomorphism
algorithm or a canonical-labelling hash.
\end{remark}

\section{Proof of the conjecture}\label{sec:proof}

\begin{proof}[Proof of \cref{thm:main}]
The case $n=10$ is \cref{prop:base10}.  Suppose $n\ge11$ and choose a
counterexample $G$ of minimum order.  The graph is non-complete and is not
$C_5$.  By \cref{lem:published10}, there are vertices $u,v$ with
$N[u]\subseteq N[v]$.

Fix any such pair and put $H=G-u$.  If $H$ is complete, then
\cref{cor:complete-deletion} gives a contradiction.  If $H$ is
$2$-connected and non-complete, there are three possibilities.  If
\[
 H\not\cong K_{n-1}^2,K_{n-1}^{n-3},
\]
then $H$ has order $n-1\ge10$ and is not one of the three exceptional graphs
for its order.  Minimality of $G$ gives
\[
  \eta(H)\ge2(n-1),
\]
and \cref{lem:published7} gives $\eta(G)\ge2n$, a contradiction.  If
$H\cong K_{n-1}^2$ or $K_{n-1}^{n-3}$, the contradiction follows from
\cref{prop:k2extension,prop:kn3extension}.

Consequently, for every dominated pair in $G$, the deletion of the dominated
vertex is not $2$-connected.  The hypotheses of \cref{lem:published9} are
therefore satisfied.  Let $C_1,\dots,C_k$ be the blocks of $G-u$ and form the
pieces $G_i$ as in \cref{sec:blocks}.  Each $G_i$ is a proper $2$-connected subgraph of $G$, and
\cref{lem:pieces-nonspecial} shows that it is not one of the three exceptional
graphs for its order.

If $s_i\ge10$, minimality therefore gives $\eta(G_i)\ge2s_i$; if $s_i<10$,
\cref{prop:pieces} gives the lower bound $L(s_i)$.  Hence
\cref{lem:blockassembly,lem:deficits} apply.  For $n\ge12$ they immediately
give $\eta(G)\ge2n$.  If $n=11$, the only unresolved size pattern is
$(6,7)$, which is settled by \cref{prop:glue}.  This final contradiction
proves the theorem.
\end{proof}

\section{Certificate and reproducibility}\label{sec:repro}

The supplementary archive contains three C++17 verifiers and the fixed
small-graph snapshots used in \cref{sec:finite}.  The command
\texttt{verify.sh} checks the SHA256 digests of the inputs, compiles the
verifiers, and reproduces the counts and minima stated in
\cref{prop:pieces,prop:glue,prop:base10}.  The fixed files are
\nolinkurl{connected7_graph6.txt}, \nolinkurl{connected7_masks.txt}, and
\nolinkurl{atlas_bases_4_7.txt}.  An optional Python script regenerates them
from NetworkX~3.6.1; NetworkX is not needed once the snapshots are present.
The C++ programs use only integer arithmetic and the algorithms described in
\cref{rem:certificate}.

\begin{thebibliography}{9}

\bibitem{BKLPS}
M.~Bonamy, M.~Knor, B.~Lu\v{z}ar, A.~Pinlou and R.~\v{S}krekovski,
\emph{On the difference between the Szeged and the Wiener index},
Appl. Math. Comput. \textbf{312} (2017), 202--213.
\href{https://doi.org/10.1016/j.amc.2017.05.047}{doi:10.1016/j.amc.2017.05.047}.

\bibitem{ReadWilson}
R.~C.~Read and R.~J.~Wilson,
\emph{An Atlas of Graphs}, Oxford University Press, Oxford, 1998.

\bibitem{NetworkXAtlas}
NetworkX Developers,
\emph{Graph generators: Atlas}, NetworkX 3.6.1 documentation,
\href{https://networkx.org/documentation/stable/reference/generated/networkx.generators.atlas.graph_atlas_g.html}{online documentation},
accessed 8 September 2026.

\end{thebibliography}

\end{document}
